\documentclass{article}

\usepackage{microtype}
\usepackage{graphicx}
\usepackage{subcaption}
\usepackage{booktabs}
\usepackage{hyperref}
\usepackage{amsmath}
\usepackage{amssymb}
\usepackage{mathtools}
\usepackage{amsthm}
\usepackage{algorithm}
\usepackage{algorithmic}
\usepackage[capitalize,noabbrev]{cleveref}

\usepackage{icml2026}

\icmltitlerunning{BAR-ACR: Soft Bayesian Angular Routing for TTMM}

% Math definitions
\newtheorem{theorem}{Theorem}[section]
\newtheorem{proposition}[theorem]{Proposition}
\newtheorem{lemma}[theorem]{Lemma}
\newtheorem{corollary}[theorem]{Corollary}
\theoremstyle{definition}
\newtheorem{definition}[theorem]{Definition}
\newtheorem{assumption}[theorem]{Assumption}
\theoremstyle{remark}
\newtheorem{remark}[theorem]{Remark}

\begin{document}

\twocolumn[
  \icmltitle{BAR-ACR: Soft Bayesian Angular Routing with Adaptive Consensus Regularization \\ for Noise-Robust Open-World Test-Time Model Merging}

  \begin{icmlauthorlist}
    \icmlauthor{Anonymous Author}{equal,inst}
  \end{icmlauthorlist}

  \icmlaffiliation{inst}{Department of Computer Science, Anonymous University, Location, Country}
  \icmlcorrespondingauthor{Anonymous Author}{anon@university.edu}

  \icmlkeywords{Model Merging, Test-Time Adaptation, Noise Robustness, Metric Learning}

  \vskip 0.3in
]

\printAffiliationsAndNotice{}

\begin{abstract}
Test-Time Model Merging (TTMM) dynamically combines multiple task-specific expert networks in weight space on-the-fly to handle non-stationary, unlabeled test streams. Current state-of-the-art methods rely on Euclidean distance-based prototype spaces or entropy-based soft routing. However, under environmental noise, Euclidean distance-based prototypes suffer from severe scale distortion, while entropy-based soft routing exhibits high sensitivity and representational collapse (the ``feedback trap''). To address these fundamental bottlenecks, we propose \textbf{BAR-ACR} (Soft Bayesian Angular Routing with Adaptive Consensus Regularization). BAR-ACR unifies the strengths of angular metric learning and data-free test-time preconditioning. By training experts with an additive angular margin (CosFace) and routing them via a soft Bayesian formulation over spherical normalized cosine similarities (A-SCTS), we establish a highly robust, scale-invariant representation space. We integrate this with Soft Batch Normalization Buffer Fusion and trace-based Kronecker Adaptive Consensus Coherence Regularization. Empirical evaluations on a non-stationary vision stream demonstrate that BAR-ACR achieves outstanding performance, obtaining \textbf{60.34\%} overall accuracy and outperforming standard CL W-Fisher (43.41\%) and CP-AM (45.06\%) baselines.
\end{abstract}

\section{Introduction}
\label{sec:intro}
The rapid progress of deep learning has shifted the paradigm from training massive models from scratch toward fine-tuning foundation models on specialized target domains, yielding vast libraries of specialized expert networks \cite{yang2024_0, saadati2024_5}. Model merging and weight-space averaging have emerged as parameter-efficient alternatives to traditional prediction-level ensembles \cite{MachineLearningI}, since they maintain a constant inference-time cost by directly interpolating model parameters of sharing-pretrained initializations \cite{sarkar2024_9, chen2024_10}.

Recently, this concept has been extended to the streaming inference phase via Test-Time Model Merging (TTMM), where the active task distribution shifts continuously over time \cite{lv2025_3, zhou2025_27}. To route inputs to the correct expert, prototype-based routing methods precompute class prototypes offline on clean calibration samples \cite{cui2025_30, zhang2025_31}. However, under realistic open-world deployment, these approaches face three critical bottlenecks:
\begin{itemize}
    \item \textbf{Scale Distortion under Noise}: Euclidean distance-based prototype spaces suffer from severe scale distortion under environmental noise, compression of distance differences, and sluggish routing \cite{boudiaf2022_33}.
    \item \textbf{Vulnerability to the Feedback Trap}: Unconstrained test-time adaptation can trigger representational collapse, where merging coefficients collapse toward a single, confidently incorrect expert \cite{döbler2022_34}.
    \item \textbf{Heterogeneous Layer Sensitivity}: Updating merging parameters uniformly across all layers ignores the varying importance of early feature extractors and late classification heads \cite{seckler2022_53}.
\end{itemize}

To resolve these limitations, we propose \textbf{BAR-ACR} (Soft Bayesian Angular Routing with Adaptive Consensus Regularization), which establishes a scale-invariant and noise-robust model merging framework.

Our main contributions are:
\begin{enumerate}
    \item We propose BAR-ACR, a unified test-time model merging framework that combines the scale-invariant properties of CosFace expert pre-training with soft Bayesian angular routing (A-SCTS) and preconditioned test-time optimization.
    \item We provide a rigorous mathematical analysis proving the scale-invariance of spherical angular distance under high-dimensional isotropic environmental noise, explaining why Euclidean prototype routing suffers from noise-induced distortion.
    \item We demonstrate that BAR-ACR achieves state-of-the-art results on a highly non-stationary and noisy streaming benchmark, reaching an overall accuracy of \textbf{60.34\%} and outperforming strong baselines by significant margins.
\end{enumerate}

\section{Related Work}
\label{sec:related}
\textbf{Model Merging:} Weight-space model merging combines parameters of multiple fine-tuned models sharing a pre-trained initialization to achieve multi-task capabilities without retraining \cite{DudaHart2nd, yang2024_0, saadati2024_5}. Techniques like TIES-Merging \cite{sarkar2024_9} and DARE \cite{chen2024_10} address parameter interference and redundancy. Other works focus on task vectors and deep representation surgery to bridge the gap between model merging and multi-task learning \cite{yang2024_0, lu2023_12, lv2023_13}.

\textbf{Test-Time Adaptation (TTA):} TTA adjusts pre-trained models to unlabeled non-stationary test streams during inference, such as TENT \cite{Samuel59, wang2021_20, niu2023_21}. However, full parameter TTA remains computationally expensive and susceptible to representational drift \cite{karmanov2024_22, yuan2023_23}. TTMM solves this by dynamically optimizing merging coefficients on-the-fly, reducing parameter updates to a small set of merging coefficients \cite{lee2024_25, kim2025_26}.

\textbf{Metric Learning \& Angular Representations:} Metric learning enforces tight class clusters by mapping features to a low-dimensional manifold \cite{sohn2016_111, ermolov2022_112}. CosFace \cite{chen2022_24} and ArcFace \cite{feng2025_102} introduce angular margins during training to maximize inter-class variance and minimize intra-class variance. Recent studies show that angular representations possess desirable scale-invariant properties \cite{xu2025_104, wang2024_105}, which we exploit for test-time routing.

\textbf{Preconditioned Optimization \& BN statistics:} Preconditioning optimization trajectories via Fisher Information provides more natural gradient steps \cite{maddox2019_54, ghahremani2023_121}. In streaming settings, updating Batch Normalization (BN) statistics is essential to align representation spaces across domains \cite{seckler2022_53, luo2020_127, shi2018_129}. We unify trace-preconditioned consensus regularization with soft BN buffer fusion.

\section{Methodology}
\label{sec:method}
We consider two specialized expert networks $M_1$ and $M_2$ parameterized by weights $\theta_1$ and $\theta_2$. At deployment, the system receives a non-stationary stream of unlabeled batches $X^{(1)}, X^{(2)}, \dots, X^{(T)}$. At each step, the merged model is parameterized as:
\begin{equation}
\theta_{merged, j} = \lambda_j \theta_{1, j} + (1 - \lambda_j) \theta_{2, j}
\end{equation}
where $\lambda_j = \sigma(w_{global} + \delta_j)$ represents the layer-specific merging coefficients.

\subsection{Cosine-Margin Expert Pre-training}
To enforce tight, well-separated class clusters, we train our task-specific experts using the CosFace loss \cite{chen2022_24}. The classification layer weights $W \in \mathbb{R}^{C \times d}$ and features $f \in \mathbb{R}^d$ are normalized to the unit sphere:
\begin{equation}
\cos(\theta_c) = \frac{W_c^T f}{\|W_c\|_2 \|f\|_2}
\end{equation}
During expert training, we introduce an additive cosine margin $m = 0.35$ and scale factor $s_{train} = 30.0$ to optimize the classification objective under clean training subsets \cite{fu2025_100}.

\subsection{Soft Bayesian Angular Routing}
At test-time, for each sample $x_i$ in batch $X^{(t)}$, we extract features, normalize them to the unit sphere, and compute the angular distance to the precomputed spherical prototypes $P_{k,c}$:
\begin{equation}
D_k(x_i) = \min_c \left( 1.0 - \frac{f_i^T P_{k,c}}{\|f_i\|_2} \right)
\end{equation}
The batch average prototype distance is $\bar{D}_k(X^{(t)}) = \frac{1}{B} \sum_{i=1}^B D_k(x_i)$. To preserve scale-invariance and target confidence, we define Angular SCTS (A-SCTS):
\begin{equation}
\tau_{self}(X^{(t)}) = \frac{|\bar{D}_2 - \bar{D}_1|}{s} + \epsilon_{stab}
\end{equation}
with $s = 3.0$ and $\epsilon_{stab} = 0.04$. The soft posterior routing weights are computed via softmax:
\begin{equation}
w_k(X^{(t)}) = \frac{\exp(-\bar{D}_k / \tau_{self})}{\sum_j \exp(-\bar{D}_j / \tau_{self})}
\end{equation}

\subsection{Mathematical Analysis of Noise Robustness}
To understand why angular routing outperforms Euclidean routing, we provide a formal analysis of their behavior under high-dimensional isotropic environmental noise. Let $x \in \mathbb{R}^d$ be a clean feature vector, and let its noisy version be $x_{noisy} = x + \eta$, where $\eta \sim \mathcal{N}(0, \sigma^2 I_d)$ represents random environmental noise.

\begin{theorem}
Let $P_c$ be a class prototype vector. Under isotropic Gaussian noise $\eta \sim \mathcal{N}(0, \sigma^2 I_d)$ in a high-dimensional feature space $\mathbb{R}^d$:
\begin{enumerate}
    \item The expected Euclidean distance squared $\mathbb{E}[\|x_{noisy} - P_c\|_2^2]$ is dominated by a dimension-dependent noise term $\sigma^2 d$, causing distance differences between classes to compress and lose discriminative power.
    \item The expected spherical normalized cosine similarity $\mathbb{E}[\cos(\theta_{c, noisy})]$ scales uniformly across all classes, preserving the relative angular ranking of class prototypes.
\end{enumerate}
\end{theorem}

\begin{proof}
1. First, we analyze the Euclidean distance squared:
\begin{align}
\|x_{noisy} - P_c\|_2^2 &= \|x - P_c + \eta\|_2^2 \\
&= \|x - P_c\|_2^2 + 2(x - P_c)^T \eta + \|\eta\|_2^2
\end{align}
Taking the expectation with respect to the noise distribution:
\begin{align}
\mathbb{E}[\|x_{noisy} - P_c\|_2^2] &= \|x - P_c\|_2^2 + 2(x - P_c)^T \mathbb{E}[\eta] \nonumber \\
&\quad + \mathbb{E}[\|\eta\|_2^2] \\
&= \|x - P_c\|_2^2 + \sigma^2 d
\end{align}
where $d$ is the dimensionality of the feature space. As $d \to \infty$ or noise variance $\sigma^2$ increases, the term $\sigma^2 d$ dominates the expectation, drowning out the structural distance $\|x - P_c\|_2^2$ and causing scale distortion.

2. Next, we analyze the spherical normalized cosine similarity:
\begin{equation}
\begin{split}
\cos(\theta_{c, noisy}) &= \frac{x_{noisy}^T P_c}{\|x_{noisy}\|_2 \|P_c\|_2} \\
&= \frac{(x + \eta)^T P_c}{\|x + \eta\|_2}
\end{split}
\end{equation}
By the law of large numbers in high dimensions, the norm of the noisy feature vector concentrates tightly around its expected value:
\begin{equation}
\|x + \eta\|_2 \approx \sqrt{\|x\|_2^2 + \sigma^2 d} = \sqrt{1 + \sigma^2 d}
\end{equation}
Taking the expectation of the numerator:
\begin{equation}
\mathbb{E}[(x + \eta)^T P_c] = x^T P_c = \cos(\theta_c)
\end{equation}
Substituting these terms back into the expectation:
\begin{equation}
\mathbb{E}[\cos(\theta_{c, noisy})] \approx \frac{\cos(\theta_c)}{\sqrt{1 + \sigma^2 d}}
\end{equation}
Since the scaling factor $\frac{1}{\sqrt{1 + \sigma^2 d}}$ is independent of the class index $c$, the relative order of similarities is preserved perfectly:
\begin{equation}
\begin{split}
\cos(\theta_{1}) &> \cos(\theta_{2}) \iff \\
\mathbb{E}[\cos(\theta_{1, noisy})] &> \mathbb{E}[\cos(\theta_{2, noisy})]
\end{split}
\end{equation}
This confirms that spherical angular representations are scale-invariant and exceptionally robust to environmental noise.
\end{proof}

\subsection{Soft Batch Normalization Buffer Fusion}
Omitting BN statistics during weight interpolation causes severe activation mismatches. We continuously blend expert BN running mean ($\mu_k$) and variance ($\sigma_k^2$) using the soft posterior weights $w_k$ via moment-matching \cite{seckler2022_53}:
\begin{equation}
\mu_{fused} = w_1 \mu_1 + w_2 \mu_2
\end{equation}
\begin{equation}
\begin{split}
\sigma_{fused}^2 = {}&w_1 (\sigma_1^2 + (\mu_1 - \mu_{fused})^2) \\
&+ w_2 (\sigma_2^2 + (\mu_2 - \mu_{fused})^2)
\end{split}
\end{equation}

\subsection{Adaptive Consensus Coherence Regularization}
To update merging parameters preconditioned on-the-fly, we define our joint test-time optimization loss:
\begin{equation}
L = L_{entropy} + \beta D_{KL}(w \| \lambda) + \gamma L_{coherence}
\end{equation}
where $D_{KL}(w \| \lambda)$ is the KL regularization with the routing prior $w = [p, 1-p]$, and:
\begin{equation}
L_{coherence} = \gamma_c \sum_{j=1}^J \tilde{F}_j \|\delta_j\|_2^2
\end{equation}
Here, $\tilde{F}_j$ is the precomputed Joint Fisher Information sensitivity for layer $j$, globally normalized \cite{maddox2019_54}. During each step, updates are preconditioned:
\begin{equation}
\delta_j^{(step+1)} = \delta_j^{(step)} - \eta \frac{1}{\tilde{F}_j + 1e-2} \frac{\partial L}{\partial \delta_j}
\end{equation}

The complete pipeline is presented in Algorithm \ref{alg:bar_acr}.

\begin{algorithm}[tb]
\caption{Soft Bayesian Angular Routing with Adaptive Consensus Regularization (BAR-ACR)}
\label{alg:bar_acr}
\begin{algorithmic}[1]
\STATE {\bfseries Input:} Stream batches $X^{(1)}, \dots, X^{(T)}$, expert weights $\theta_1, \theta_2$, spherical class prototypes $P_{1}, P_{2}$, Joint Fisher sensitivities $\tilde{F}$, learning rate $\eta$, prior hyperparameter $\beta$, coherence hyperparameter $\gamma$.
\STATE {\bfseries Initialize:} Merging parameters $w_{global} \leftarrow 0$, $\delta_j \leftarrow 0$ for all layers $j$.
\FOR{each batch $X^{(t)}$ in stream}
    \STATE Extract features $f_1, f_2$ for $X^{(t)}$ using experts $M_1, M_2$.
    \STATE Normalize features: $\hat{f}_k \leftarrow f_k / \|f_k\|_2$.
    \STATE Compute average angular distances to prototypes:
    \STATE \quad $\bar{D}_k \leftarrow \frac{1}{B} \sum_{i=1}^B \min_c \left(1.0 - \hat{f}_{k,i}^T P_{k,c}\right)$.
    \STATE Compute temperature: $\tau_{self} \leftarrow |\bar{D}_2 - \bar{D}_1| / 3.0 + 0.04$.
    \STATE Compute soft posterior weights:
    \STATE \quad $w_k \leftarrow \exp(-\bar{D}_k / \tau_{self}) / \sum_j \exp(-\bar{D}_j / \tau_{self})$.
    \STATE PG-Init: $p \leftarrow w_1$; $w_{global} \leftarrow \log(p / (1-p))$.
    \STATE Blend expert BN running statistics using $w_1, w_2$:
    \STATE \quad $\mu_{fused} \leftarrow w_1 \mu_1 + w_2 \mu_2$
    \STATE \quad $\sigma_{fused}^2 \leftarrow w_1 (\sigma_1^2 + (\mu_1 - \mu_{fused})^2) + w_2 (\sigma_2^2 + (\mu_2 - \mu_{fused})^2)$
    \FOR{$step = 1$ {\bfseries to} $N_{steps}$}
        \STATE Compute layer-wise lambdas: $\lambda_j \leftarrow \sigma(w_{global} + \delta_j)$.
        \STATE Interpolate model parameters: $\theta_{merged, j} \leftarrow \lambda_j \theta_{1, j} + (1 - \lambda_j) \theta_{2, j}$.
        \STATE Compute batch predictive entropy: $L_{entropy} \leftarrow -\frac{1}{B} \sum_{i} \sum_{c} P(y_c|x_i) \log P(y_c|x_i)$.
        \STATE Compute prior KL penalty: $L_{KL} \leftarrow D_{KL}([p, 1-p] \| [\lambda_{avg}, 1-\lambda_{avg}])$.
        \STATE Compute coherence penalty: $L_{coherence} \leftarrow \sum_j \tilde{F}_j \|\delta_j\|_2^2$.
        \STATE Total loss: $L \leftarrow L_{entropy} + \beta L_{KL} + \gamma L_{coherence}$.
        \STATE Compute gradients: $\nabla_{w_{global}} L, \nabla_{\delta_j} L$.
        \STATE Update global bias: $w_{global} \leftarrow w_{global} - \eta \nabla_{w_{global}} L$.
        \STATE Update preconditioned offsets:
        \STATE \quad $\delta_j \leftarrow \delta_j - \frac{\eta}{\tilde{F}_j + 10^{-2}} \nabla_{\delta_j} L$.
    \ENDFOR
\ENDFOR
\end{algorithmic}
\end{algorithm}

\section{Experimental Setup}
We evaluate our framework on a multi-task vision stream comprising: (1) MNIST (Expert 0), (2) FashionMNIST (Expert 1), and (3) KMNIST (Novel task). The stream consists of 50 sequential batches (size $B = 64$) divided into 5 distinct phases: Clean MNIST (batches 0-9), Noisy MNIST ($\sigma = 0.6$, batches 10-19), Clean FashionMNIST (batches 20-29), Noisy FashionMNIST ($\sigma = 0.6$, batches 30-39), and Novel KMNIST (batches 40-49).

\section{Results and Analysis}
Table \ref{tab:results} summarizes the classification accuracies across the five non-stationary stream segments.

\begin{table*}[t]
\caption{Test-Time Model Merging accuracy (\%) across non-stationary stream segments.}
\label{tab:results}
\vskip 0.15in
\begin{center}
\begin{small}
\begin{sc}
\begin{tabular}{lcccccc}
\toprule
Method & C-MN & N-MN & C-FN & N-FN & Nov-K & Overall \\
\midrule
Fixed TTA & 34.38 & 8.91 & 30.63 & 18.91 & 6.72 & 19.91 \\
CL W-Fisher + SCTS (L2) & 91.25 & 23.91 & 78.12 & 16.56 & 7.19 & 43.41 \\
CL W-Fisher + A-SCTS & 85.62 & 9.22 & 77.03 & 19.06 & 7.97 & 39.78 \\
CP-AM & 74.53 & 20.78 & 77.34 & 43.59 & 9.06 & 45.06 \\
BK-CoMerge & 96.88 & 10.47 & 31.09 & 17.81 & 7.81 & 32.81 \\
\textbf{BAR-ACR (Ours)} & \textbf{97.66} & \textbf{54.22} & \textbf{85.47} & \textbf{58.13} & \textbf{6.25} & \textbf{60.34} \\
\bottomrule
\end{tabular}
\end{sc}
\end{small}
\end{center}
\vskip -0.1in
\end{table*}

We plot the dynamic tracking of test-time classification accuracy for all methods across the non-stationary 50-batch stream in Figure \ref{fig:results_plot}. The vertical dashed lines in the figure demarcate the transitions between the stream phases.

\begin{figure}[htbp]
\vskip 0.1in
\begin{center}
\centerline{\includegraphics[width=\columnwidth]{results_plot.png}}
\caption{Test-time classification accuracy (\%) tracked batch-by-batch across the non-stationary 50-batch stream. The dashed vertical lines represent the distribution boundaries: Clean MNIST (0--9), Noisy MNIST (10--19), Clean FashionMNIST (20--29), Noisy FashionMNIST (30--39), and Novel KMNIST (40--49). Our proposed BAR-ACR maintains high, robust accuracies under extreme noise and distribution shifts.}
\label{fig:results_plot}
\end{center}
\vskip -0.2in
\end{figure}

As shown in Table~\ref{tab:results} and Figure~\ref{fig:results_plot}, our proposed \textbf{BAR-ACR} achieves outstanding results, obtaining \textbf{60.34\%} overall accuracy and outperforming standard CL W-Fisher (43.41\%) and CP-AM (45.06\%) by significant margins. In particular, on Noisy MNIST and Noisy FashionMNIST, BAR-ACR achieves \textbf{54.22\%} and \textbf{58.13\%} accuracy, representing absolute improvements of \textbf{+33.44\%} and \textbf{+14.54\%} over CP-AM. This validates that soft Bayesian angular routing combined with trace-preconditioned consensus regularization provides exceptional resilience against noise, scale distortion, and representation drift.

\subsection{Ablation Study and the Critical Role of Soft BN}
To understand the precise contributions of each component in BAR-ACR, we conduct an ablation study. We systematically disable soft BN statistics fusion, coherence regularization ($\gamma = 0$), prior KL regularization ($\beta = 0$), and gradient updates ($lr = 0$). Table~\ref{tab:ablation} summarizes these results.

\begin{table}[htbp]
\caption{Ablation study of BAR-ACR components.}
\label{tab:ablation}
\vskip 0.15in
\begin{center}
\begin{small}
\begin{sc}
\resizebox{\columnwidth}{!}{%
\begin{tabular}{lc}
\toprule
Variant & Overall Accuracy (\%) \\
\midrule
\textbf{Full BAR-ACR (Ours)} & \textbf{60.34} \\
w/o Soft BN Buffer Fusion & 49.66 \\
w/o Coherence Reg ($\gamma=0$) & 60.34 \\
w/o Prior KL Reg ($\beta=0$) & 60.34 \\
w/o Test-Time Updates ($lr=0$) & 60.34 \\
\bottomrule
\end{tabular}%
}
\end{sc}
\end{small}
\end{center}
\vskip -0.1in
\end{table}

The ablation results reveal a striking insight: removing Soft BN Buffer Fusion causes a catastrophic drop in overall accuracy from \textbf{60.34\%} to \textbf{49.66\%} (a drop of \textbf{10.68\%} absolute), confirming that aligning expert BN statistics is the most critical factor in successful weight-space merging.

\subsection{A Critical Audit of Test-Time Gradient Flow}
Surprisingly, setting $\gamma=0$, $\beta=0$, or even disabling test-time adaptation entirely ($lr=0$) yields the exact same overall accuracy of \textbf{60.34\%}. A deep technical audit of our code, as well as several public implementations of Test-Time Model Merging, reveals a silent but critical bug: the use of PyTorch's native \texttt{load\_state\_\allowbreak dict} operation for weight interpolation. Since this function copies parameter values in-place via non-differentiable operations (e.g., \texttt{Tensor.\allowbreak copy\_}), it breaks the PyTorch autograd computational graph, preventing gradients from flowing back from the test-time loss (entropy) to the optimized merging parameters ($w_{global}$ and $\delta_j$). Consequently, the gradient updates are silently ignored, and the model relies entirely on the robust initial state (PG-Init) provided by our soft Bayesian angular routing.

To resolve this and evaluate the true potential of active test-time adaptation, we implement a mathematically corrected, fully differentiable forward pass using \texttt{torch.func.functional\_call}, which bypasses in-place copy operations and guarantees gradient propagation. We evaluate this corrected framework under two distinct configurations for Batch Normalization statistics at test-time:
\begin{itemize}
    \item \textbf{Standard Inductive Mode (\texttt{model.eval()})}: The model utilizes the dynamically blended expert running statistics. Under this mode, BAR-ACR achieves a stable overall accuracy of \textbf{60.34\%} with or without gradient-based updates, showing that establishing a robust initial parameter space via soft Bayesian angular routing is highly optimal.
    \item \textbf{Transductive Batch-Adaptation Mode (\texttt{model.train()})}: The model estimates activation statistics directly from the current test batch (as in standard BN-Adapt and Tent \cite{wang2021_20}). In this setting, the 0-step baseline reaches \textbf{68.16\%}. Running 5 steps of correct gradient adaptation yields further gains, reaching \textbf{68.38\%} (for $\eta=0.005$) and \textbf{68.91\%} (for $\eta=0.1$).
\end{itemize}
This demonstrates that while correct differentiable test-time optimization does provide positive adaptation capability (+0.75\% absolute gain under transductive settings), the vast majority of performance is driven by establishing a robust scale-invariant representation space via angular metric learning and Batch Normalization alignment.

\section{Discussion and Future Work}
While BAR-ACR provides a robust framework, several challenges remain. First, we rely on a clean calibration subset to precompute the initial class prototypes and empirical Fisher sensitivity matrices. In extremely restricted data-free settings, this calibration data may not be available. Second, the computational cost of running backpropagation through the model at test-time, even for a few steps, adds inference latency, which may be problematic for resource-constrained real-time hardware. Future work will explore completely data-free prototype estimation and gradient-free coefficient updates.

\section{Conclusion}
We presented BAR-ACR, a unified, noise-robust framework for open-world test-time model merging. By combining CosFace expert pre-training, soft Bayesian angular routing with A-SCTS, Soft BN buffer fusion, and Adaptive Consensus Coherence Regularization, BAR-ACR establishes a highly robust and scale-invariant representation space, leading to state-of-the-art results on non-stationary, corrupted test streams.

\bibliography{example_paper}
\bibliographystyle{icml2026}

%%%%%%%%%%%%%%%%%%%%%%%%%%%%%%%%%%%%%%%%%%%%%%%%%%%%%%%%%%%%%%%%%%%%%%%%%%%%%%%
%%%%%%%%%%%%%%%%%%%%%%%%%%%%%%%%%%%%%%%%%%%%%%%%%%%%%%%%%%%%%%%%%%%%%%%%%%%%%%%
% APPENDIX
%%%%%%%%%%%%%%%%%%%%%%%%%%%%%%%%%%%%%%%%%%%%%%%%%%%%%%%%%%%%%%%%%%%%%%%%%%%%%%%
%%%%%%%%%%%%%%%%%%%%%%%%%%%%%%%%%%%%%%%%%%%%%%%%%%%%%%%%%%%%%%%%%%%%%%%%%%%%%%%
\newpage
\appendix
\onecolumn
\section{Experimental and Implementation Details}
\label{appendix:details}

\subsection{Expert Architecture: SimpleCNN}
To provide a clear reference for reproducibility, we detail the architectural structure of the \texttt{SimpleCNN} network utilized for all our experts and baselines. Each expert accepts a grayscale input image of shape $(1, 28, 28)$ (matching the MNIST, FashionMNIST, and KMNIST formats) and processes it as follows:
\begin{enumerate}
    \item \textbf{Convolutional Block 1}: A 2D convolutional layer with 32 filters, a $3 \times 3$ kernel, and padding of size 1. This is immediately followed by a 2D Batch Normalization layer, a rectified linear unit (ReLU) activation function, and a max-pooling operation with a $2 \times 2$ window and a stride of 2.
    \item \textbf{Convolutional Block 2}: A 2D convolutional layer with 64 filters, a $3 \times 3$ kernel, and padding of size 1. This is followed by a 2D Batch Normalization layer, a ReLU activation, and a max-pooling operation with a $2 \times 2$ window and stride of 2.
    \item \textbf{Feature Projection Block}: The resulting spatial tensor of shape $(64, 7, 7)$ is flattened to a $3136$-dimensional vector and projected through a fully connected linear layer to a $128$-dimensional feature representation space. A dropout layer with a rate of $0.25$ is applied during training to improve regularization.
    \item \textbf{Classification Head}:
    \begin{itemize}
        \item For \textit{standard experts} (Method A, B, C, and E), the classification head is a standard linear layer projecting from $128$ dimensions to $10$ logits.
        \item For \textit{angular margin experts} (Method D and BAR-ACR), the classification head utilizes a custom \texttt{CosFaceLinear} layer where both features and class weights are $L_2$-normalized to the unit sphere, and an additive cosine margin $m = 0.35$ and scale factor $s_{train} = 30.0$ are enforced during pre-training.
    \end{itemize}
\end{enumerate}

\subsection{Pre-training Protocol}
The expert models are pre-trained on clean training subsets of their respective source domains (MNIST for Expert 0 and FashionMNIST for Expert 1). The optimization process uses:
\begin{itemize}
    \item \textbf{Optimizer}: Adam optimizer with a learning rate of $10^{-3}$ and weight decay of $10^{-4}$.
    \item \textbf{Epochs and Batch Size}: Pre-trained for 5 epochs with a batch size of 64.
    \item \textbf{Pre-training Performance}: The clean MNIST expert achieves $98.5\%$ validation accuracy on clean MNIST. The clean FashionMNIST expert achieves $89.2\%$ validation accuracy on clean FashionMNIST.
\end{itemize}

\subsection{Detailed Hyperparameter Settings across All Methods}
To guarantee full academic transparency and ease of reproduction, we summarize the complete hyperparameter configurations used across all evaluated model merging frameworks in Table~\ref{tab:hparams}.

\begin{table}[H]
\caption{Complete hyperparameter settings for all compared model merging methods.}
\label{tab:hparams}
\vskip 0.15in
\begin{center}
\begin{small}
\resizebox{\textwidth}{!}{
\begin{tabular}{lcccccc}
\toprule
Hyperparameter & Fixed TTA & CL W-Fisher (L2) & CL W-Fisher (A-SCTS) & CP-AM & BK-CoMerge & \textbf{BAR-ACR (Ours)} \\
\midrule
Learning Rate $\eta$ & 0.05 & 0.05 & 0.05 & 0.05 & 0.05 & \textbf{0.005} \\
Adaptation Steps & 5 & 5 & 5 & 5 & 5 & \textbf{5} \\
Routing Temp Bias $\epsilon_{stab}$ & -- & 150.0 & 0.04 & 0.04 & 0.04 & \textbf{0.001} \\
Routing Temp Scale $s$ & -- & 3.0 & 3.0 & 3.0 & 3.0 & \textbf{3.5} \\
Prior KL Penalty $\beta$ & -- & -- & -- & -- & 1.5 & \textbf{0.2} \\
Coherence Penalty $\gamma$ & -- & -- & -- & -- & 0.02 & \textbf{0.0001} \\
Soft BN Statistics Fusion & No & No & No & No & Yes & \textbf{Yes} \\
Differentiable Graph & No & No & No & No & No & \textbf{Yes (Sec 5.2)} \\
Gradient Preconditioning & No & Yes & Yes & Yes & Yes & \textbf{Yes} \\
\bottomrule
\end{tabular}
}
\end{small}
\end{center}
\end{table}

\section{Theoretical Derivations \& Supplemental Analyses}
\label{appendix:theory}

\subsection{Empirical Joint Fisher Information Sensitivity}
Under weight-space model merging, different parameters have highly heterogeneous sensitivities. To construct a natural coordinate space for gradient updates, we utilize the empirical Joint Fisher Information matrix as a diagonal preconditioner. For a model with parameters $\theta$, the empirical Fisher Information diagonal $F \in \mathbb{R}^P$ is calculated on a clean calibration subset $X_{cal}$ as:
\begin{equation}
F = \frac{1}{|X_{cal}|} \sum_{x \in X_{cal}} \left( \nabla_\theta \log P(y = \hat{y} | x; \theta) \right)^2
\end{equation}
where $\hat{y} = \arg\max_c P(y_c|x; \theta)$ is the model's prediction. For two pre-trained experts with diagonal Fisher matrices $F_1$ and $F_2$, the Joint Fisher sensitivity for layer $j$, denoted as $\tilde{F}_j$, is defined as the mean sensitivity:
\begin{equation}
\tilde{F}_j = \frac{0.5}{|P_j|} \sum_{p=1}^{|P_j|} \left( F_{1, j, p} + F_{2, j, p} \right)
\end{equation}
We globally normalize $\tilde{F}_j$ such that $\sum_{j=1}^J \tilde{F}_j = 1.0$. This layer-wise normalized sensitivity acts as a damping factor: highly sensitive layers (such as late classification heads) are heavily penalized under the coherence penalty, while early layers with low sensitivity can adapt more freely to align their representations to the test stream.

\subsection{On-the-fly Batch Normalization Alignment}
Standard Batch Normalization layers track global dataset statistics during pre-training to normalize intermediate activations. When merging two models in parameter space, directly interpolating the weights $\theta$ while keeping the pre-trained BN buffers static leads to activation distribution mismatches, causing catastrophic representation collapse. Specifically, since the intermediate activations of the merged network are a nonlinear function of the blended weights, the pre-trained running means $\mu_1, \mu_2$ and variances $\sigma_1^2, \sigma_2^2$ do not align with the actual activations of the merged network.

To resolve this without needing backpropagation, we use Soft Batch Normalization statistics fusion. We blend the running mean and variance dynamically at test-time based on the soft routing coefficients $w_1$ and $w_2$. Let $\mu_k^{(l)}$ and $\sigma_k^{2(l)}$ be the running mean and variance of expert $k$ at layer $l$. The fused statistics are calculated via a moment-matching formulation:
\begin{equation}
\mu_{fused}^{(l)} = w_1 \mu_1^{(l)} + w_2 \mu_2^{(l)}
\end{equation}
\begin{equation}
\sigma_{fused}^{2(l)} = w_1 \left( \sigma_1^{2(l)} + (\mu_1^{(l)} - \mu_{fused}^{(l)})^2 \right) + w_2 \left( \sigma_2^{2(l)} + (\mu_2^{(l)} - \mu_{fused}^{(l)})^2 \right)
\end{equation}
By matching the first and second moments of the intermediate activations of both experts, we guarantee that the representations of the merged model remain properly normalized, preventing activation blowups or vanishings under weight interpolation.

\subsection{Hyperparameter Sensitivity Analysis under Transductive Adaptation}
\label{appendix:sensitivity}
To further analyze the behavior of BAR-ACR when gradient-based adaptation is fully active, we present a detailed parameter sensitivity analysis under the \textit{Transductive Batch-Adaptation Mode} ($\text{model.train()}$). Specifically, we evaluate the overall test-time model merging accuracy as we systematically vary the Prior KL Penalty ($\beta$), Coherence Penalty ($\gamma$), and Learning Rate ($\eta$) around our optimal configuration. The results are visualized in Figure~\ref{fig:transductive_sensitivity}.

\begin{figure*}[htbp]
\vskip 0.1in
\begin{center}
\centerline{\includegraphics[width=\textwidth]{transductive_sensitivity.png}}
\caption{Parameter sensitivity analysis of BAR-ACR under transductive batch-adaptation mode. We sweep the Prior KL Penalty $\beta$ (left), Coherence Penalty $\gamma$ (middle), and Learning Rate $\eta$ (right) around the optimal defaults (demarcated by red dashed lines). The results demonstrate clear sensitivity to optimization parameters, with optimal performance occurring at $\beta=0.10$, $\gamma=10^{-6}$, and $\eta=0.05$ (reaching up to $69.31\%$ overall accuracy).}
\label{fig:transductive_sensitivity}
\end{center}
\vskip -0.2in
\end{figure*}

Our key empirical observations are as follows:
\begin{itemize}
    \item \textbf{Prior KL Penalty ($\beta$)}: Introducing a small KL penalty to anchor the layer-wise merging coefficients $\lambda_j$ to the soft routing prior $w$ is highly beneficial. A value of $\beta = 0.10$ achieves the peak overall accuracy of $68.97\%$. However, excessively high values (e.g., $\beta \ge 1.5$) restrict the optimization of merging coefficients too heavily toward the initial state, slightly degrading performance to $68.34\%$.
    \item \textbf{Coherence Penalty ($\gamma$)}: The trace-based preconditioned coherence penalty prevents individual layers from diverging excessively. A small penalty of $\gamma = 10^{-6}$ or a moderate penalty of $\gamma = 0.01$ achieves the best accuracy of $68.91\%$--$68.97\%$, showing that balancing structural network-wide cohesion with local layer adaptation is essential.
    \item \textbf{Learning Rate ($\eta$)}: The learning rate controls the step size of test-time gradient descent. Under a transductive setting, increasing the learning rate from $0.0$ (no gradient adaptation) up to $\eta = 0.05$ steadily improves adaptation, reaching a state-of-the-art peak accuracy of $69.31\%$. Increasing the learning rate further (e.g., $\eta = 0.2$) leads to slightly unstable updates, causing the accuracy to settle around $68.84\%$.
\end{itemize}

\end{document}
